\documentclass[12pt,twoside,fleqn,a4paper]{article}
\makeatletter %%%% ugly fudge to customise spacing in lists %%%%%%%%%%%%%%%
  %% note \@listI is originally defined in art12.sty (or size12.clo for 2e)
  \def\@listI{\leftmargin\leftmargini  \parsep\z@ plus2.5\p@ minus\p@
    \topsep\z@ plus4\p@ minus6\p@  \itemsep 2\p@ plus2.5\p@ minus\p@}
  \let\@listi\@listI
  %% could define \@listii, \@listiii similarly for sublists.
\makeatother

\usepackage{graphicx}

\pagestyle{plain}
\setlength{\textwidth}{160mm}
\setlength{\textheight}{245mm}
\setlength{\parskip}{1.4ex}
\setlength{\parindent}{0cm}
\setlength{\oddsidemargin}{0cm}
\setlength{\evensidemargin}{0cm}
\setlength{\topmargin}{-0.5cm}
\setlength{\headsep}{0.5cm}
\setlength{\headheight}{1cm}

\raggedbottom

\setlength{\partopsep}{-0.5\parskip}%
\setlength{\leftmargini}{2em}%
\setlength{\leftmarginii}{2em}%
\setlength{\labelsep}{0.5em}%
\setlength{\topsep}{0\parskip}%

% Define italic ``et al''. Use as {\etal} to avoid space problems
\newcommand{\etal}{{\it et al.}}

\begin{document}

\section{Introduction}

We look at the gap-filling problem in metabolic processes.

Here, we have a set of metabolites, and a set of reactions. In steady
state, each reaction uses up metabolites at a given rate, and produces
metabolites at given rate.

We have a set of experiments where we have measured the metabolites
used, and the metabolites produced. We know (or think we know) some of
the reactions involved. We wish to work out the other reactions
involved. Each potential reaction has a notional 'cost' reflecting the
probability it is involved; e.g. this reaction has been seen in
similar contexts, or this reaction is from a plant database, so is
unlikely to be used.

The decision variables are the \emph{flux}, or throughput of each
potential reaction. The constraints reflect the mass balance
equations: the amount used and produced must balance with the observed
use and output of the experiments.

\section{Nomenclature}

We have metabolites in the set $M = \{1 .. m\}$, and reactions in the
set $R = \{1 .. r\}$.

The \emph{flux} $v_i$ determines how active is reaction $i$. $v_i$ is
the principal decision variable of the problem. The 'cost' of reaction
$i$ is $c_i$.

Each reaction involves a number of metabolites. The stoichiometric
coefficient $a_{ik}$ determines the amount of metabolite $k$ used or
produced in reaction $i$: a negative coefficient means the metabolite
is used, a positive coefficient means the metabolite is
produced. There is a linear relationship between the flux and the
amounts: for negative stoichiometric coefficients, the total amount of
metabolite $k$ used by reaction $i$ is $v_i . -a_{ik}$ (since $a_{ik}$
will be negative). Similarly the amount of metabolite produced is
$v_i . a_{ik}$ (since $a_{ik}$ will be positive).

The rate of use of source metabolites -- that is, metabolites drawn
from the substrate, and not supplied by reactions -- is given by $s_k$.
The rate of production of output metabolites is given by $p_k$.

\begin{tabular}{c | c c c c c c  | c c}
        & R1 &  R2 &  R3 &  R4 &  R5 &  Biomass & & \\ \hline
 Cost   & $c_1$ & $c_2$ & $c_3$ & $c_4$ & $c_5$ & -1000 & \\ \hline
 Flux   & $v_1$ & $v_2$ & $v_3$ & $v_4$ & $v_5$ & $v_6$ & LB & UB \\ \hline
  Met 1 & -1 &   0 &   1 &   1 &   0 &   0 & -10 & 0\\
  Met 2 &  1 &  -1 &   0 &   1 &   0 &  -1 & 0 & 0  \\
  Met 3 &  0 &  -1 &   0 &   0 &  -1 &   0 & 0 & 0   \\
  Met 4 &  0 &   0 &  -1 &  -1 &   0 &   0 & 0 & 0   \\
  Met 5 &  0 &   1 &   0 &  -1 &   0 &   1 & 0 & 10   \\
  Met 6 &  0 &   0 &   0 &  -1 &   1 &   0 & 0 & 0
\end{tabular}

\section{Formulation 1}

The problem can be expressed simply:

Problem {\bf P1}:

\begin{eqnarray}
  \mbox{minimise} \sum_i c_i v_i & & 
\end{eqnarray}
subject to
\begin{eqnarray}
  \sum_i v_i a_{ik} & = & p_k - s_k \;\;\; \forall k \in M \label{eqn:metbal}
\end{eqnarray}

Note, we add dummy reactions in $D$ to $R$, each of which produce or consume
one metabolite. These ensure that the problem is feasible,
allowing an arbitrary amount of any metabolite to either appear or
disappear if required. The cost for use of these reactions is, of
course, set very high.

Note also that one chosen reaction will be a biomass production
reaction. This reaction is kind-of the aim of the process - to produce
biomass. The biomass production reaction can have an negative
objective coefficient so that it is maximised rather than
minimised. The absolute value of the objective coefficient needs to be
chosen carefully so that the biomass achieves maximum production given
supply restrictions.

\section{Dual variables}

We can solve this problem using dual variables. We start by solving
problem {\bf P1} where only dummy reactions are available. We
then examine dual costs of every reaction not currently included. This
is a simple calculation involving only the cost of the reaction, and the
dual variables of each of the metabolites either used or produced by
the reaction (along with the stoichiometric coefficients).

Let $\lambda_k$ be the dual variable associated with the constraint
(\ref{eqn:metbal}) on metabolite $k$. The reduced cost of reaction $i$
is
\[
\mbox{RC}_i = c_i - \sum_k{\lambda_k a_{ik}} \label{eqn:rc1}
\]

We add any reactions with a negative reduced cost to the the LP.  We
can the solve the new LP, and repeat until no reactions have negative
reduced cost.

\section{Integer Formulation}

The ``GlobalFit'' system uses an integer formulation, using variables
$y_i$ = 1 if reaction $i$ is used; 0 otherwise. We use an upper bound
on flux ($\mbox{UB}_i$) for each reaction. The forumlation is then

Problem {\bf P2}:

\begin{eqnarray}
  \mbox{minimise} \sum_i c_i y_i & & 
\end{eqnarray}
subject to
\begin{eqnarray}
  \sum_i v_i a_{ik} & = & p_k - s_k \;\;\; \forall k \in M \label{eqn:metbal2} \\
  v_i & \le & \mbox{UB}_i \,  y_i \;\;\; \forall i \in R \\
  y_i & \in & \{0, 1\} \;\;\; \forall i \in R 
\end{eqnarray}

\section{Reachability}

The formulations in the previous section do not consider the order of
production of metabolites.

For instance, consider a reaction $r$ that consumes metabolite $A$, and
produces metabolite $B$. Another reaction $s$ consumes metabolites $B$, and
produces metabolite $A$. Neither $A$ nor $B$ are available in the
substrate (i.e. $s_A = s_B = 0$). In that case, neither reaction can
be used, unless another source of $A$ or $B$ is available. But
Formulation 1 will happily use \emph{only} $r$ and $s$ since the use and
production of $A$ and $B$ balance out.

This gives rise to the concept of \emph{reachability}, and
\emph{depth}.

We say a reaction \emph{enabled} if all the metabolites consumed by a reaction
are available. We say a metabolite is available if it is either
supplied in the substrate ($s_k > 0$), or a reaction the supplies it
is enabled.

We define the 'level' of a reaction as the level it is enabled in the
following algorithm. Note: \emph{never} is a constant indicating that
the reaction can not be enabled.

The \emph{depth} of the residuals can also be calculated: when all
metabolites with $p_k$ > 0 are \emph{available}, we know the residuals
can theoretically be produced.


Algorithm \emph{Depth}:

\begin{itemize}
\item Make all metabolites $k$ with $s_k > 0$ available
\item Set \emph{depth} of all reactions to {\bf never}
\item Set \emph{residual\_depth} to {\bf never}
\item \emph{level} = 0
\item \emph{finished} = {\bf false}
\item \emph{is\_progress } = {\bf true}
\item while {\bf not} \emph{finished} {\bf and} \emph{is\_progress}
  \begin{itemize}
  \item \emph{level} = \emph{level + 1}
  \item \emph{finished} = {\bf true}
  \item \emph{is\_progress} = {\bf false}
  \item for each reaction $r$ not yet enabled
    \begin{itemize}
    \item if all metabolites required by $r$ are available
      \begin{itemize}
      \item enable reaction $r$
      \item set \emph{depth} of $r$ to \emph{level}
      \item \emph{is\_progress} = {\bf true}
      \end{itemize}
    \item else
      \begin{itemize}
      \item \emph{finished} = {\bf false}
      \end{itemize}
    \end{itemize}
  \item for each reaction $r$ enabled at this level
    \begin{itemize}
    \item for each metabolite $i$ produced by reaction $r$
      \begin{itemize}
      \item make $i$ available
      \end{itemize}
    \end{itemize}
  \item if all metabolites $k$ with $p_k > 0$ are available
    \begin{itemize}
    \item All residuals can now be produced
    \item set \emph{residual\_depth} = \emph{level}
    \item set \emph{finished} = {\bf true}
    \end{itemize}
  \end{itemize}
\end{itemize}

The variable \emph{is\_progress} ensures there is progress at each
iteration, or else the algorithm terminates. The variable
\emph{finished} is left at {\bf true} when all reactions are enabled.

\section{Depth is important}

Depth can be used to help ensure that the solution can be
implemented in real life.

We note the there must be at least one one reaction from each level in
an implementable solution. Once the substrate metabolites are made
available, only reactions in level 1 can be triggered. Hence, at least
one of those must be used to produce metabolites that are not
otherwise available. The same is true at level 2: I can't use any of
those reactions until I see the metabolites produced at level 1. But
if I want to use reactions in level 3, I must use at least one of the
reactions in level 2. This repeats until I reach the depth of the
residual metabolites.

This then gives rise to a ``valid inequality'': at least one equation
from each depth must be used.

\section{Formulation 2}

Let $d$ be the residual-depth as calculated by algorithm \emph{Depth}.

Define $J(r) = j$ for $j \in \{1 .. d\}$ define the level of reaction 
$r \in R$. Define \emph{LB} as a lower bound on a
flux for a used reaction. Then we can define Formulation 2 as
Formulation 1 plus

\begin{eqnarray}
  \sum_{i : J(i) = j} v_i & \ge & LB \;\;\; \forall j = 1..d \label{eqn:depth}
\end{eqnarray}


\section{Dual variables revisited}

Let $\mu_j$ be the dual variable associated with the constraint
(\ref{eqn:depth}) on reactions in depth $j$. The reduced cost of reaction
$i$ is now
\[
\mbox{RC}_i = c_i - \sum_k{\lambda_k a_{ik}} - \mu_j \label{eqn:rc2}
\]

\section{Compartments}

Cells are modelled to have ``compartments'' - different areas where
reactions take place. Typical compartments are:
\begin{itemize}
\item Cytoplasm
\item Periplasm
\item External
\item Outside
\end{itemize}

External is external to the cell, and the source of normal supply
(i.e. the substrate we are growing on). We also model an ``outside''
compartment, which is also external, but is the site for biomass
production and other {\it residuals} from processes.

We simply expand all reactions and metabolites to be present
in/operate in a compartment. Transport reactions move metabolites from
one compartment to another (e.g. $CO_2$ in Cytoplasm can be moved to
$CO_2$ in the Outside.).

\section{Multiple experiments}

Rather than having a single supply/residual figure, biologists will
run several experiments. Having measured outcomes with one substrate
(i.e. supply configuration), they may knock out one metabolite from
the substrate, and run the experiment again to see if biomass is
created.

In this case, the same reaction can be used with different fluxes in
each experiment. A different reaction can be used in different
experiments, but we want to find the smallest set of reactions that
works in all cases. So, if I can explain 2 experiments with a slightly
higher flux through 2 reactions, rather than a smaller flux through 3
reactions, I want to choose the solution with 2 reactions.

Expand our nomenclature to include $e$ experiments, and the set $E =
\{1..e\}$ of all experiments.

Data is supply/production data (including biomass production) $s^j_k$ and
$p^j_k$ on $j \in E$.

We now need flux variables $v^j_i$, for $j \in E$. We
still use an overall flux of $v_i$, and ensure that $v_i$ takes the
maximum value amongst $v^j_i$ values.

Problem {\bf P2}:

\begin{eqnarray}
  \mbox{minimise} \sum_i c_i v_i & & 
\end{eqnarray}
subject to
\begin{eqnarray}
  \sum_i v^j_i a_{ik} & = & p^j_k - s^j_k \;\;\; \forall k \in
  M, j \in E \label{eqn:metbal3} \\
  v^j_i & \le & v_i \;\;\; \forall i \in R, j \in E\label{eqn:expflux}
\end{eqnarray}


\end{document}

% LocalWords:  stoichiometric ik Reachability reachability Periplasm
